% Preámbulo compartido para todas las clases de Beamer.
% Cada clase lo incluye con % Preámbulo compartido para todas las clases de Beamer.
% Cada clase lo incluye con % Preámbulo compartido para todas las clases de Beamer.
% Cada clase lo incluye con % Preámbulo compartido para todas las clases de Beamer.
% Cada clase lo incluye con \input{../../latex/preamble.tex}

\documentclass[aspectratio=169,11pt]{beamer}

\usetheme{Madrid}
\setbeamertemplate{navigation symbols}{}
\setbeamertemplate{footline}[frame number]

\usepackage[utf8]{inputenc}
\usepackage[spanish]{babel}
\usepackage{amsmath,amssymb}
\usepackage{graphicx}
\usepackage{booktabs}
\usepackage{listings}
\usepackage{xcolor}

\lstset{
  basicstyle=\ttfamily\small,
  keywordstyle=\color{blue},
  commentstyle=\color{gray},
  stringstyle=\color{teal},
  showstringspaces=false,
  breaklines=true,
  frame=single
}

\newcommand{\definicion}[1]{%
  \begin{block}{Definición}
    #1
  \end{block}
}

\newcommand{\ejemplo}[1]{%
  \begin{exampleblock}{Ejemplo}
    #1
  \end{exampleblock}
}

\newcommand{\importante}[1]{%
  \begin{alertblock}{Importante}
    #1
  \end{alertblock}
}

\author{Agustín Gurvich}
\institute{Introducción a Machine Learning}


\documentclass[aspectratio=169,11pt]{beamer}

\usetheme{Madrid}
\setbeamertemplate{navigation symbols}{}
\setbeamertemplate{footline}[frame number]

\usepackage[utf8]{inputenc}
\usepackage[spanish]{babel}
\usepackage{amsmath,amssymb}
\usepackage{graphicx}
\usepackage{booktabs}
\usepackage{listings}
\usepackage{xcolor}

\lstset{
  basicstyle=\ttfamily\small,
  keywordstyle=\color{blue},
  commentstyle=\color{gray},
  stringstyle=\color{teal},
  showstringspaces=false,
  breaklines=true,
  frame=single
}

\newcommand{\definicion}[1]{%
  \begin{block}{Definición}
    #1
  \end{block}
}

\newcommand{\ejemplo}[1]{%
  \begin{exampleblock}{Ejemplo}
    #1
  \end{exampleblock}
}

\newcommand{\importante}[1]{%
  \begin{alertblock}{Importante}
    #1
  \end{alertblock}
}

\author{Agustín Gurvich}
\institute{Introducción a Machine Learning}


\documentclass[aspectratio=169,11pt]{beamer}

\usetheme{Madrid}
\setbeamertemplate{navigation symbols}{}
\setbeamertemplate{footline}[frame number]

\usepackage[utf8]{inputenc}
\usepackage[spanish]{babel}
\usepackage{amsmath,amssymb}
\usepackage{graphicx}
\usepackage{booktabs}
\usepackage{listings}
\usepackage{xcolor}

\lstset{
  basicstyle=\ttfamily\small,
  keywordstyle=\color{blue},
  commentstyle=\color{gray},
  stringstyle=\color{teal},
  showstringspaces=false,
  breaklines=true,
  frame=single
}

\newcommand{\definicion}[1]{%
  \begin{block}{Definición}
    #1
  \end{block}
}

\newcommand{\ejemplo}[1]{%
  \begin{exampleblock}{Ejemplo}
    #1
  \end{exampleblock}
}

\newcommand{\importante}[1]{%
  \begin{alertblock}{Importante}
    #1
  \end{alertblock}
}

\author{Agustín Gurvich}
\institute{Introducción a Machine Learning}


\documentclass[aspectratio=169,11pt]{beamer}

\usetheme{Madrid}
\setbeamertemplate{navigation symbols}{}
\setbeamertemplate{footline}[frame number]

\usepackage[utf8]{inputenc}
\usepackage[spanish]{babel}
\usepackage{amsmath,amssymb}
\usepackage{graphicx}
\usepackage{booktabs}
\usepackage{listings}
\usepackage{xcolor}

\lstset{
  basicstyle=\ttfamily\small,
  keywordstyle=\color{blue},
  commentstyle=\color{gray},
  stringstyle=\color{teal},
  showstringspaces=false,
  breaklines=true,
  frame=single
}

\newcommand{\definicion}[1]{%
  \begin{block}{Definición}
    #1
  \end{block}
}

\newcommand{\ejemplo}[1]{%
  \begin{exampleblock}{Ejemplo}
    #1
  \end{exampleblock}
}

\newcommand{\importante}[1]{%
  \begin{alertblock}{Importante}
    #1
  \end{alertblock}
}

\author{Agustín Gurvich}
\institute{Introducción a Machine Learning}
